\section{Introduction}\label{sec:introduction}

Let $X$ be a smooth integral complex projective variety of dimension $n\ge3$,
and let $L$ be an ample line bundle. Put $V=L^n$. For a tuple
$\mathbf p=(p_1,\ldots,p_r)$ of distinct points, the ordinary multipoint
Seshadri constant is
\begin{equation}\label{eq:seshadri-definition}
 \eps(X,L;\mathbf p)
 =\inf_C\frac{L\cdot C}{\sum_{i=1}^r\mult_{p_i}C},
\end{equation}
where $C$ ranges over integral curves meeting at least one of the marked
points, and the multiplicity is zero at points outside $C$. This invariant
measures the local positivity of $L$ at the entire tuple. Its universal
upper bound is
\begin{equation}\label{eq:volume-bound}
 \eps(X,L;\mathbf p)\le (V/r)^{1/n}.
\end{equation}
Equality means that no curve imposes a stronger constraint than the volume
of $L$.

Write
\[
 U_r(X)=\{(p_1,\ldots,p_r)\in X^r:p_i\ne p_j\text{ when }i\ne j\}.
\]
A property holds at \emph{very general} tuples if the exceptional tuples
are contained in a countable union of proper Zariski-closed subsets of
$U_r(X)$.

\begin{theorem}\label{thm:main}
Let $X$ be a smooth integral complex projective variety of dimension
$n\ge3$, and let $L$ be an ample line bundle on $X$. There is an integer
$r_0=r_0(X,L)\ge1$ such that for every integer $r\ge r_0$ there is a
countable union $Z_r$ of proper Zariski-closed subsets of $U_r(X)$, with
nonempty complement, for which
\[
 \eps(X,L;\mathbf p)=\left(\frac{L^n}{r}\right)^{1/n}
 \qquad\text{for every }\mathbf p\in U_r(X)\setminus Z_r.
\]
\end{theorem}

\begin{samepage}
Equivalently, on the blow-up $\pi:Y\to X$ at such a tuple, with exceptional
divisors $E_1,\ldots,E_r$, the real divisor class
\[
 \pi^*L-(V/r)^{1/n}\sum_{i=1}^rE_i
\]
is nef and has top self-intersection zero. Here nefness means nonnegative
intersection with every integral curve. The threshold depends on the fixed
polarized variety, but all integers beyond it are included. In particular,
Theorem~\ref{thm:main} asserts exact equality for each such point count,
rather than convergence of a normalized constant as $r$ increases.
\end{samepage}

\subsection{Background and relation to earlier work}

The plane problem originates in Nagata's work on Hilbert's fourteenth
problem. He conjectured a lower bound on the degree of a plane curve in
terms of its multiplicities at very general points, and proved the strict
inequality when the number of points is a square at least $16$
\cite{Nagata1959}. Multipoint Seshadri constants express the corresponding
positivity question on an arbitrary polarized variety. The qualitative
Nagata--Biran--Szemberg formulation asks for maximality at every sufficiently
large point count. Ro\'e and Ross state this formulation in arbitrary
dimension and distinguish the ordinary curve-defined constant from its
higher-cycle analogues \cite[Conjecture~1.3 and Remark~1.4]{RoeRoss2009}.
We use the ordinary constant throughout, even when $\dim X>2$.

Biran's packing stability theorem supplies an important symplectic
precedent: sufficiently many equal balls can fill a closed symplectic
four-manifold with rational symplectic class with arbitrarily small volume
deficit \cite{Biran1999}; see also \cite[Theorem~6.3]{Biran2001}. The algebraic
question fixes the complex variety and asks for a nef class on its point
blow-ups, so symplectic packing stability is not by itself the required
algebraic statement. In arbitrary dimension, K\"uchle obtained lower bounds
for multipoint Seshadri constants that are asymptotically optimal as the
number of points tends to infinity \cite[Theorem~1.1]{Kuechle1996}.
Biran's product inequality, extended to higher-cycle constants by
Ro\'e--Ross, relates point counts on a variety to point counts on projective
space \cite[Theorem~1.1 and Remark~1.2]{RoeRoss2009}.
These results illuminate the volume scale in \eqref{eq:volume-bound}; an
asymptotically optimal estimate alone does not imply eventual equality.

The companion paper \cite{OpenAI2026Surfaces} establishes the qualitative
surface statement by a nodal exponent bound, diagonal slice compression,
and a primitive lattice direction tuned to the number of points. Its final
tests are \emph{injective} jet evaluations on finite torus subgroups.
We adapt the planar geometric mechanism to polynomial spaces in two
variables and use \emph{surjective} full-dimensional jet evaluations instead.
All arguments required for this adaptation are proved here; the surface
maximality theorem is not an input.

The interpretation of positivity through asymptotic jet generation is
classical; see Demailly \cite[Section~6]{Demailly1992}.
Our use of initial monomials belongs to the valuation approach to linear
series initiated by Okounkov \cite{Okounkov1996} and developed by Lazarsfeld--Musta\c{t}\u{a}
\cite{LazarsfeldMustata2009} and Kaveh--Khovanskii
\cite{KavehKhovanskii2012}. More directly, Ito compares multipoint jet
separation with initial monomial spaces by separate finite-degree
degenerations, without a finite-generation hypothesis
\cite[Proposition~5.7 and Theorem~5.8]{Ito2013}.
Monomial diagrams and interpolation minors also underlie Dumnicki's
diagram-cutting method \cite[Proposition~13 and Theorem~14]{Dumnicki2007}.
We retain elementary proofs of the specific transfer and counting statements
used below.

\subsection{Proof strategy and the two-coordinate construction}

Separating jets through degree $m$ at a tuple means prescribing,
independently at each point, all Taylor coefficients of total degree at
most $m$, in local coordinates and line-bundle frames. If $H^0(X,kL)$
separates these jets, then
\[
 \eps(X,L;\mathbf p)\ge \frac{m}{k};
\]
the curve-intersection argument is given in Lemma~\ref{lem:jets-curves}.
Thus, for $w=(V/r)^{1/n}$, our interpolation target is jet separation
through degree $\lfloor k\theta\rfloor$ for every fixed $0<\theta<w$
and all sufficiently large $k$. We first construct a tuple for each such
test and later obtain simultaneous tests at very general tuples.

The argument works with finite exponent sets
$S_k\subset\Z^n$, one for each positive integer degree $k$, satisfying
\[
 S_k+S_l\subset S_{k+l}.
\]
For a set $S$, its monomial space is the span of the Laurent monomials
$z^\alpha$ with $\alpha\in S$ on $(\C^*)^n$. The multiplication condition
will be as important as the geometry of the individual sets.

A complete-intersection flag on $X$ first supplies sets $S_k$ whose
cardinalities are $h^0(X,kL)$. They lie in the dilates of a simplex
\[
 \Delta(a_1,\ldots,a_n)
 =\left\{x\in\R_{\ge0}^n:\sum_{i=1}^n\frac{x_i}{a_i}\le1\right\}
\]
of volume $V/n!$. Jet surjectivity for the corresponding monomials can be
transferred back to sections on $X$.

The main construction replaces any two intercepts $A,B$ of the bounding
simplex by $w,AB/w$, for any sufficiently small positive real $w$.
It preserves the number of exponents, the multiplication condition, and
the implication from jet surjectivity after the replacement to jet
surjectivity before it. The other exponent coordinates are carried along
unchanged. The planar node bound and primitive-direction argument come
from the surface method; the real-intercept replacement and its compatibility
with these remaining coordinates allow the process to be iterated in
higher dimensions.

For a fixed sufficiently large integer $r$, choose $w=(V/r)^{1/n}$.
Repeated replacements produce sets $F_k$ inside dilates of
\[
 P_* = \Delta(w,\ldots,w,rw),
 \qquad \#F_k=h^0(X,kL),
 \qquad F_k+F_l\subset F_{k+l}.
\]
Because this simplex still has volume $V/n!$, only $o(k^n)$ of its degree-$k$
lattice points are missing. The multiplication condition upgrades this
density statement: for each fixed compact set $K\subset\operatorname{int}P_*$,
every lattice point of $kK$ belongs to $F_k$ for all sufficiently large $k$.
Indeed, such a point has on the order of
$k^n$ decompositions into two nearly equal degrees, whereas the holes can
exclude only $o(k^n)$ of them. This short argument is a full-density version
of the interior approximation results for semigroups
\cite[Theorem~1.6]{KavehKhovanskii2012}.

The long last intercept of $P_*$ permits explicit interpolation at $r$
points that differ only in their last coordinate. One-variable Hermite
interpolation, applied to the coefficients of monomials in the other
coordinates, gives arbitrary jets through degree $\lfloor k\theta\rfloor$
for every fixed $0<\theta<w$ and all sufficiently large $k$.
Transferring these tests to $X$, taking a countable intersection of
nonempty Zariski-open loci, and then letting $\theta$ increase to $w$
proves Theorem~\ref{thm:main}.

Section~\ref{sec:transfers} establishes the monomial operations and their
jet transfers, and Section~\ref{sec:flag} constructs the initial exponent
sets on $X$. Section~\ref{sec:reshaping} proves the two-coordinate
replacement. Section~\ref{sec:saturation} establishes interior filling and
interpolation, and Section~\ref{sec:conclusion} assembles the very-general
conclusion.
